\section*{The Propers: Detailed Commentary}

\subsection*{Beholding, truth, and the Christian name\enspace\elementref{Entrance; Collect}}

Psalm 17's entrance clause comes from a prayer that contrasts integrity of path with violent opposition. Its desire to behold and be satisfied is therefore not aesthetic curiosity. The worshipper seeks vindication and final communion before God. As an adapted Missal antiphon it gives a compact gathering horizon.

Augustine's exposition of the psalm (on the Vulgate/LXX Psalm 16) treats the speaker as Christ together with his Body and reads the final satisfaction eschatologically: the hunger of exile ends when divine glory is manifested. Augustine leaves the psalm's historical questions open, and his reading guards the antiphon from becoming a promise of immediate vindication. The Entrance teaches desire to wait for God's face rather than demanding that enemies, suffering, or ambiguity disappear before worship can begin.

The Collect holds mercy and moral discernment together. God shows truth's light to those who wander and recalls them; the petition then asks Christians to reject what conflicts with their name and pursue what fits it. “Christian” is neither a label that sanctifies any conduct nor an achievement earned before grace. Identity is received in Christ and becomes a norm for conversion.

The verbs also prevent truth and mercy from being separated into pastoral alternatives. Light is shown to persons who have wandered precisely so that return is possible; return then includes both renunciation and pursuit. The prayer neither humiliates wanderers as a permanent class nor treats recall as affirmation of every path. It asks the whole assembled Church, not an absent group, to submit its Christian name to continuing judgment and grace.

\subsection*{Rain, word, and an earth genuinely affected\enspace\elementref{First Reading; Responsorial Psalm}}

Isaiah 55 compares two effective descents. Rain and snow do not return before watering, fertilizing, and providing seed and bread; God's word likewise does not return empty but accomplishes the purpose for which it is sent. The analogy asserts divine efficacy without making hearers inert. In the chapter, the same voice commands thirsting people to come, the wicked to forsake their way, and the hearer to return for abundant pardon.

The Syriac witness Aphrahat quotes Isaiah 55:10--11 in \emph{Demonstration VIII}, first to insist that the divine word actually completes its mission and later within an argument about resurrection. His Christological and eschatological reuse is not a continuous exposition of Isaiah's exilic oracle, but it shows how early reception heard more than successful information-transfer: the word descends, acts in embodied history, and returns with its work accomplished. In its literal sense the oracle first addresses displaced Judah; Christian reception reads the same promise on a wider horizon.

Aphrahat's two uses also disclose an interpretive discipline. In section 6 the oracle establishes divine reliability; in section 15 it participates in a network of resurrection promises. The same verse can therefore bear more than one ecclesial use without being made infinitely plastic. Its controlling feature remains the sent word's effective return. That feature links promise, Incarnation, proclamation, and final fulfillment.

Psalm 65 supplies a liturgical answer in the world of creatures. God visits land, fills watercourses, breaks clods, softens earth, and blesses yield. Human sowing is not denied, but the Psalm's praise begins with divine visitation. The response, marked Luke 8:8, overlays the image of fruitful soil; that marker does not turn Luke into the Mass's Gospel or authorize replacing Matthew's distinct form.

Augustine's exposition of the Vulgate/LXX Psalm 64 follows the psalm from pardon and worship into watered furrows and harvest. He reads the rain as truth given in a form hearers can receive and warns that fruit gives no ground for boasting in unaided merit. He reads the psalm allegorically rather than in its historical sense, and his reading maintains the order of gift, cultivation, and fruit. The soil is genuinely worked and genuinely fruitful, yet water and increase remain received.

Cyril of Alexandria's \emph{Sermon XLI on Luke} directly expounds the source of the response. He identifies Christ as the sower, distinguishes path, rock, thorns, and fertile ground, and observes that Matthew's thirty-, sixty-, and hundredfold yields signify unequal measures among the saints rather than unequal value before God. Because the Lectionary response is Lukan while the Gospel is Matthean, Cyril receives the refrain's own source, and Luke's explanation does not replace Matthew's proclaimed wording.

\Needspace{2\baselineskip}
This reception also clarifies the refrain's work between Isaiah and Paul. The hundredfold response voices confidence that divine gift can become abundant creaturely fruit, but the verses sung beneath it still praise water, earth, and harvest rather than classifying hearers. The refrain anticipates the Gospel by echo, while the Psalm retains its own grammar of communal thanksgiving. Cyril's reception deepens the overlay, and both textual layers stay audible.

The liturgical transition is therefore cumulative without making the sources interchangeable.

The pairing can support ecological gratitude without reducing either text to environmental policy. Land, water, seed, bread, flocks, and human praise belong together under gift. Yet Romans' groaning creation warns against pretending that present fertility is universal or that crop failure marks deficient faith.

\subsection*{Creation's labor and the redemption of bodies\enspace\elementref{Second Reading}}

Romans 8:18--23 begins with comparison, not denial. Present sufferings are real; coming glory exceeds them. Creation is personified as straining forward, subjected to futility not by its own choice and yet held in hope. Its promised freedom is not annihilation but release from corruption into the freedom associated with God's children.

Believers share the tension. They possess the Spirit's firstfruits and still groan; they have adoption and still await its bodily completion. Chrysostom reads creation's labor toward incorruption and emphasizes that the liberated body remains God's creature. Hope therefore does not despise matter, romanticize pain, or claim that every temporal loss will be visibly repaired before resurrection.

In \emph{Homily 14 on Romans}, Chrysostom attends to Paul's rhetorical personification rather than treating creation as a decorative metaphor. Creation is made to speak because its renewal displays the generosity shown to humanity; corruption is not its final state, and the believers' firstfruits intensify rather than cancel waiting. His reading keeps three subjects distinct---creation, believers, and the Spirit whose treatment follows after the appointed boundary---and so resists collapsing every groan into one vague spirituality.

Gregory of Nyssa invokes Romans 8:19--23 in \emph{Against Eunomius} IV while arguing about Christ as firstborn. He imagines the angelic creation sharing the grief of human perdition and awaiting the manifestation of restored children. Gregory reuses the passage for doctrine rather than defining Paul's “creation,” and his reading differs from Chrysostom's emphasis on the nonrational cosmos clothed with incorruption. The two readings together are instructive: ancient interpreters agreed that redemption is not a solitary interior event, but they did not delimit the personified creation identically.

The difference also checks easy ecological appropriation. Chrysostom strongly supports the future of nonhuman creation; Gregory widens cosmic solidarity through the angelic world. Neither supplies a modern environmental program. The appointed text grounds hope for created reality and bodily liberation, from which practical responsibilities may be reasoned only through further historical and moral steps.

The passage is read in semi-continuous course. “Creation,” “earth,” and “fruit” may resonate with Isaiah, the Psalm, and Gospel. It belongs first to Pauline eschatology.

\subsection*{A composed acclamation and two genuinely different proclamation paths\enspace\elementref{Acclamation; long Gospel; short Gospel}}

The acclamation is a liturgical composition. The official U.S. page gives it no biblical citation. Its theological role is still clear: it welcomes Christ about to speak and names him as sower of the word. That identification belongs to the appointed liturgical text even though it is not presented as a verbatim Bible verse.

Both Gospel forms proclaim the same parable. The sower casts broadly; three failures and one fruitful outcome expose not defects in the seed but different receptions. Path, shallow rock, and thorns portray loss of understanding, rootlessness under pressure, and the choking force of anxiety and wealth only in Jesus' interpretation within the long form. Good soil hears, understands, and bears fruit in differentiated abundance.

Chrysostom's \emph{Homily 44 on Matthew} directly treats the public parable. He calls Christ the sower, doctrine the seed, and human souls the land, stressing the sower's unsparing breadth: the same good seed is not rationed according to a prior judgment about who looks promising. Failure therefore cannot be blamed upon a defective word, while the act of sowing itself refuses despair about apparently resistant ground. \emph{Homily 45} then follows the long form into the disciples' question, Isaiah's hard-hearing oracle, and the differing yields. Chrysostom treats thirty, sixty, and a hundred as real fruit at different measures, not as one quota below which grace is absent.

Cyril's parallel Lukan sermon adds a counterweight. He applies the soils vigorously to hardheartedness, untested enthusiasm, persecution, and the growth of worldly anxiety, but concludes by urging sterile ground to awake and become cultivated. Cyril's moral diagnosis therefore names responses, not fixed classes of persons. Both Cyril and Chrysostom preach movement, perseverance, and cultivation; neither gives a modern minister warrant to infer a person's entire spiritual condition from disability, education, poverty, trauma, or one visible failure.

Origen offers a further, deliberately indirect reception in \emph{On First Principles} III.1.17. Answering an objection about unequal timing in divine instruction, he compares God to a husbandman who may delay sowing rocky ground until the seed can take secure root rather than spring up superficially. Origen does not comment continuously on Matthew 13 here; he reuses its agrarian logic to defend patient divine pedagogy. The proposal is more speculative than Chrysostom's homily, yet it supplies a valuable pastoral possibility: delay need not mean abandonment, and apparently resistant ground need not be declared incapable of future cultivation.

Jesus' parable does not say this of its failed seed; the possibility comes from Origen, not from the literal plot. Its force is negative and hopeful: one historical response is not a final divine verdict, and rapid religious enthusiasm is not always deeper than slower formation.

The long form adds a difficult discourse about parables, hearing, and Isaiah 6. Jesus' words are judicial and medicinal: hardened perception is exposed, yet conversion and healing remain the stated good refused. The beatitude of seeing and hearing names gift, not intellectual superiority.

The short form stops at “Whoever has ears ought to hear.” Its congregation does not hear Jesus identify each soil, though the canonical continuation still informs the parable. The branch instead leaves the parable's summons more open and may call for especially careful explanation of it.

\begin{studybox}{Soils are not castes}
The parable diagnoses response; it does not authorize fixed spiritual castes. A person can move from distraction to attention, from shallow enthusiasm to rooted fidelity, or from anxious attachment toward freedom. Trauma, disability, poverty, or limited education must not be labeled “bad soil.” The Gospel judges obstacles and calls for conversion while grace continues to sow.
\end{studybox}

\subsection*{Gifts ordered toward sanctification\enspace\elementref{Prayer over Offerings}}

The Prayer over the Offerings addresses the gifts of the supplicating Church and asks that they advance sanctification. The petition avoids two reductions: gifts are not payment buying holiness, and interior holiness does not make the offered signs irrelevant. Ecclesial offering is taken into Christ's sacrificial action and ordered toward the offerers' transformation.

\subsection*{Dwelling or abiding at Communion\enspace\elementref{Communion A; Communion B; Prayer after Communion}}

Communion Antiphon A takes the worshipper into Psalm 84's sanctuary longing: even birds find a home near the altars, and those dwelling in God's house continue in praise. At Communion the branch can signify Eucharistic approach as ecclesial homecoming. It does not imply static enclosure; the complete pilgrimage Psalm also knows journey, strength, petition, and trust.

Augustine's exposition of the Vulgate/LXX Psalm 83 reads heart and flesh together in the sparrow and dove. The heart seeks its lasting home; the body makes a temporary nest by works of mercy; the divine house culminates in unwearied praise. His detailed bird allegory goes beyond the psalm's literal meaning, yet it prevents “dwelling” from becoming passive enclosure: embodied charity belongs to the pilgrimage toward the house where desire is fulfilled.

Alternative B appoints John 6:57's mutual-abiding language. Augustine's direct reception of John 6 joins eating Christ with abiding in his Body and ecclesial unity. Augustine's reading bears on the Johannine antiphon only and does not turn Matthew's sower into a Eucharistic allegory in every detail.

Aquinas's \emph{Commentary on John} 6, lecture 7, distinguishes sacramental eating from its full spiritual effect without separating them. Ordinary food is transformed into the eater; Christ's food, received worthily in faith and charity, transforms the receiver toward Christ. Mutual abiding therefore names grace and incorporation rather than spatial enclosure. Aquinas also states the moral qualification: the sacrament is objectively given, yet its life-giving fruit is not described as automatic in one who resists the mode of abiding.

Augustine and Aquinas converge on ecclesial and moral participation while using different conceptual tools. Augustine emphasizes the Body's unity; Aquinas the manner and effect of reception. The unresolved branch may be illuminated by both, but neither belongs to the Psalm 84 alternative: the Missal prints no composite antiphon.

The Prayer after Communion asks that receiving gifts and frequenting the mystery deepen saving effect. Frequency is ordered repetition, not mechanism. Worthy, attentive participation can root charity across time; repetition without conversion is no guarantee of fruit. Conversely, absence of felt consolation does not prove that grace is absent.

\section*{The Propers: Notable and Quotable}

{\small
\begin{itemize}[itemsep=0em,topsep=0.15em,parsep=0pt]
\item \textbf{“A sower went out to sow” --- Earthseed after social collapse.\enspace\elementref{both Gospel forms}} Octavia E.~Butler's 1993 \emph{Parable of the Sower} turns the biblical title into dystopian science fiction: a young Black woman crosses a climate-ravaged California while cultivating a new community and the change-centered religion Earthseed. The seed is no longer only a preached word; it is a portable social experiment whose reception remains dangerously uncertain.

\item \textbf{The sower --- from Gospel notebook to blazing modern image.\enspace\elementref{both Gospel forms}} Vincent van Gogh recorded that the parable made a deep impression, copied Gospel parables in several languages, later called Christ a supreme artist of spoken parable, and returned to the sower in more than thirty visual variations. In the best-known setting-sun images, an agricultural laborer becomes a haloed figure of life, death, and artistic vocation---without becoming a mere illustration of Matthew.

\item \textbf{“The seed fell on stony ground” --- while the boy digs into limestone.\enspace\elementref{both Gospel forms}} In William Allen White's 1909 novel \emph{A Certain Rich Man}, a mother's attempts to interest young John in Beecher's sermons fail because bees, birds, and butterflies hold his attention, “so the seed fell on stony ground.” The next sentences place him above a limestone ledge, digging into a spring: spiritual idiom and literal geology quietly collide.

\item \textbf{“Creation is groaning” --- climate-policy pressure.\enspace\elementref{Second Reading}} The Lutheran World Federation's climate-justice statement says creation groans under the weight of human action and inaction, then presses for a binding international climate agreement. Paul's woman-in-labor image becomes institutional advocacy: the groan is treated not as passive doom but as evidence demanding political action.
\end{itemize}
}

\section*{The Propers: Interpretive Possibilities}

\noindent\emph{The following exploratory proposals were developed in an AI-assisted editorial process. They are attributed to none of the cited authorities and claim no historical compositional intent. Each states its anchors and limits.}

\Needspace{10\baselineskip}
\noindent\textbf{The word's journey is a circulation, not a delivery.\enspace\elementref{First Reading; Responsorial Psalm; Gospel}}
Isaiah's rain-word descends, yields seed and bread, and returns; Psalm 65 moves through water, harvest, and praise; Matthew's seed multiplies. Together they form a gift--reception--transformation--return circuit: speech becomes nourishment, new seed, and answer. Separate reading misses this circulation. Its fruit is teaching as cultivation enabling further giving, not delivery. Isaiah and Matthew officially correlate and the Psalm responds; the circuit is editorial, and no checked source, including the near analogues located, develops it. It neither mechanizes grace nor erases Matthew's losses and differing yields.

\Needspace{10\baselineskip}
\noindent\textbf{Creation becomes a hearer with more than one voice.\enspace\elementref{Responsorial Psalm; Second Reading; Gospel}}
Matthew makes soil hear; Psalm 65 makes creation sing; Romans gives it groaning labor and waiting. The movement from watered object to choir to hopeful subject enlarges hearing beyond the mind. Separate reading misses its bodily and cosmic conditions. Its fruit is ecological discipleship serving truthful receptivity under bodily redemption. Near analogues leave the conjunction undeveloped. Romans is semi-continuous, poetry does not settle creaturely consciousness, and ecology neither exhausts nor replaces redemption in Christ.

\Needspace{10\baselineskip}
\noindent\textbf{The short form lets the parable behave like seed.\enspace\elementref{Acclamation; short Gospel; long Gospel}}
The acclamation names Christ as sower. The short Gospel plants images and stops at “hear”; the long adds the parable question, Isaiah, beatitude, and interpretation. Thus image works in memory, or explanation becomes further seed still requiring perseverance. An inventory misses this changed experience. Its fruit is preaching that cultivates attention without equating explanation with reception. Chrysostom's parable awakening inquiry is a near analogue, not branch analysis. Neither form is superior; canonical context remains, and the enacted branch may not be inferred.

\Needspace{10\baselineskip}
\noindent\textbf{The bird that steals seed can also find a home at the altar.\enspace\elementref{Gospel; Communion A}}
Matthew's birds consume exposed seed; Communion A, if appointed, gives Psalm 84's birds a home near God's altars. Place reverses the image: exposure permits loss, dwelling fruitfulness. Separate reading misses movement from scattered attention to a habitat of worship and mercy. Augustine supplies a precedent: before interpreting the sanctuary birds, he warns of fruit devoured by birds. The fruit is shelter for mission, not enclosure. This does not prove Missal intent; Communion A is unresolved, the birds differ, and the sower is no Eucharistic code.

\Needspace{10\baselineskip}
\noindent\textbf{The harvest's deepest measure may be abiding rather than output.\enspace\elementref{Gospel; Communion B; Prayer after Communion}}
Matthew quantifies yields; Communion B, if chosen, names abiding through eating; the final prayer asks participation to deepen saving effect. Conformity with the gift becomes fruit's source. Separate reading misses this critique of metrics. Its fruit is ministry freed from numerical triumphalism without excusing sterility. Aquinas's food transforming the worthy receiver is a near analogue. Communion B is unresolved; abiding does not erase fruit, and frequency is neither automatic nor separable from faith, charity, and conversion.
